\documentclass[12pt,letterpaper]{article}
\usepackage[utf8]{inputenc}
\usepackage[spanish]{babel}
\usepackage{geometry}
\usepackage{array}
\usepackage{tabularx}
\usepackage{graphicx}

\geometry{margin=2.5cm}

\begin{document}

\begin{center}
\textbf{Universidad Nacional Abierta y a Distancia}\\
\textbf{Vicerrectoría Académica y de Investigación}\\
\textbf{Unidad gestora: Escuela de Ciencias Básicas, Tecnología e Ingeniería ECBTI}\\
\textbf{Programa: Ciencias Básicas}\\
\textbf{Curso: Cálculo Diferencial}\\
\textbf{Código: 100410}\\
\vspace{0.5cm}
\Large\textbf{Rúbrica de evaluación -- Tarea 2}
\end{center}

\section*{Tabla 1. Rúbrica de evaluación}

\begin{tabularx}{\textwidth}{|>{\bfseries}p{4cm}|X|}
\hline
\multicolumn{1}{|c|}{\textbf{Elemento}} & \multicolumn{1}{c|}{\textbf{Descripción}} \\
\hline
Resultado de aprendizaje: & Reconocer funciones elementales y sus propiedades, a través de conceptos básicos, representaciones gráficas y la resolución de ejercicios, para resolver situaciones específicas de manera adecuada. \\
\hline
Tipo de actividad: & Independiente \\
\hline
Puntaje de la actividad: & 135 puntos \\
\hline
\textbf{Primer criterio de evaluación:}

\textbf{Contenido:}
El estudiante desarrolla analíticamente 4 ejercicios sobre funciones.

Este criterio tiene una valoración máxima de: \textbf{50 puntos} & 
\textbf{Nivel alto:} Desarrolla correctamente de manera analítica los cuatro ejercicios asignados.

\textit{Si su trabajo se encuentra en este nivel puede obtener entre 40 puntos y 50 puntos}

\textbf{Nivel Medio:} Desarrolla correctamente de manera analítica entre dos y tres ejercicios asignados.

\textit{Si su trabajo se encuentra en este nivel puede obtener entre 30 puntos y 39 puntos}

\textbf{Nivel bajo:} Desarrolla correctamente de manera analítica solo un ejercicio asignado.

\textit{Si su trabajo se encuentra en este nivel puede obtener entre 1 puntos y 29 puntos}

\textbf{Nivel no presentado:} No desarrolla ni presenta los ejercicios en el espacio respectivo.

\textit{En este nivel se obtienen 0 puntos} \\
\hline
\end{tabularx}

\newpage

\begin{tabularx}{\textwidth}{|>{\bfseries}p{4cm}|X|}
\hline
\textbf{Segundo criterio de evaluación:}

\textbf{Contenido:}
El estudiante realiza la interpretación y análisis gráfico de ejercicios a través de la aplicación GeoGebra.

Este criterio tiene una valoración máxima de: \textbf{21 puntos} & 
\textbf{Nivel alto:} El estudiante realiza una correcta interpretación y análisis gráfico en la aplicación GeoGebra, en todos los ejercicios que así lo solicitan.

\textit{Si su trabajo se encuentra en este nivel puede obtener entre 17 puntos y 21 puntos}

\textbf{Nivel Medio:} El estudiante realiza una correcta interpretación y análisis gráfico en la aplicación GeoGebra en algunos ejercicios que así lo solicitan.

\textit{Si su trabajo se encuentra en este nivel puede obtener entre 13 puntos y 16 puntos}

\textbf{Nivel bajo:} El estudiante realiza interpretación y análisis gráfico en la aplicación GeoGebra en algunos de los ejercicios, pero lo realizado no es correcto.

\textit{Si su trabajo se encuentra en este nivel puede obtener entre 1 puntos y 12 puntos}

\textbf{Nivel no presentado:} No presenta la evidencia relacionada a la interpretación y análisis gráfico de los ejercicios.

\textit{En este nivel se obtienen 0 puntos} \\
\hline
\textbf{Tercer criterio de evaluación:}

\textbf{Procedimiento:}
El estudiante demuestra apropiación de los conceptos abordados, evidenciado a través de la sustentación en video del ejercicio de aplicación sobre funciones y realiza una breve presentación personal en inglés.

Este criterio tiene una valoración máxima de: \textbf{45 puntos} & 
\textbf{Nivel alto:} El estudiante sustenta de manera correcta a través de video la solución del ejercicio de aplicación, y evidencia dominio y apropiación de las temáticas abordas en la unidad 1, así mismo, realiza una breve presentación personal en inglés.

\textit{Si su trabajo se encuentra en este nivel puede obtener entre 36 puntos y 45 puntos}

\textbf{Nivel Medio:} El estudiante sustenta a través de video la solución del ejercicio de aplicación, sin embargo, esta no es correcta. No se evidencia dominio y apropiación de las temáticas abordas en la unidad 1, igualmente la presentación realizada en inglés no cumple con lo establecido en las indicaciones de la guía.

\textit{Si su trabajo se encuentra en este nivel puede obtener entre 27 puntos y 35 puntos}

\textbf{Nivel bajo:} El estudiante sustenta a través de video la solución del ejercicio de aplicación, pero no evidencia dominio alguno y apropiación de las temáticas abordas en la unidad 1, igualmente, no realiza la presentación personal en inglés.

\textit{Si su trabajo se encuentra en este nivel puede obtener entre 1 puntos y 26 puntos}

\textbf{Nivel no presentado:} No presenta la evidencia relacionada con la sustentación del ejercicio de aplicaciones de la Unidad 1.

\textit{En este nivel se obtienen 0 puntos} \\
\hline
\end{tabularx}

\newpage

\begin{tabularx}{\textwidth}{|>{\bfseries}p{4cm}|X|}
\hline
\textbf{Cuarto criterio de evaluación:}

\textbf{Formal:}
Los ejercicios son elaborados utilizando el editor de ecuaciones de Word y abordan el uso adecuado de la sintaxis matemática.

Este criterio tiene una valoración máxima de: \textbf{10 puntos} & 
\textbf{Nivel alto:} Los 5 ejercicios son elaborados utilizando el editor de ecuaciones de Word y abordan correctamente la sintaxis matemática.

\textit{Si su trabajo se encuentra en este nivel puede obtener entre 8 puntos y 10 puntos}

\textbf{Nivel Medio:} Entre 2 y 4 ejercicios son elaborados utilizando el editor de ecuaciones de Word y abordan correctamente la sintaxis matemática.

\textit{Si su trabajo se encuentra en este nivel puede obtener entre 6 puntos y 7 puntos}

\textbf{Nivel bajo:} Solo 1 ejercicio es elaborado utilizando el editor de ecuaciones de Word y no aborda correctamente la sintaxis matemática.

\textit{Si su trabajo se encuentra en este nivel puede obtener entre 1 puntos y 5 puntos}

\textbf{Nivel no presentado:} No realiza los ejercicios a través del editor de ecuaciones de Word.

\textit{En este nivel se obtienen 0 puntos} \\
\hline
\textbf{Quinto criterio de evaluación:}

\textbf{Participación:}
Interactúa de forma oportuna, adecuada y respetuosa en el foro de la actividad de forma semanal, respondiendo a los ejercicios propuestos de acuerdo con su selección.

Este criterio tiene una valoración máxima de: \textbf{9 puntos} & 
\textbf{Nivel alto:} Participa semanalmente en el foro con aportes significativos, de manera oportuna, adecuada y respetuosa.

\textit{Si su trabajo se encuentra en este nivel puede obtener entre 8 puntos y 9 puntos}

\textbf{Nivel Medio:} Participa con aportes significativos en algunas semanas, de manera oportuna, adecuada y respetuosa.

\textit{Si su trabajo se encuentra en este nivel puede obtener entre 6 puntos y 7 puntos}

\textbf{Nivel bajo:} No participa significativamente en el foro de la actividad.

\textit{Si su trabajo se encuentra en este nivel puede obtener entre 1 puntos y 5 puntos}

\textbf{Nivel no presentado:} No participa en el foro de la actividad con sus aportes.

\textit{En este nivel se obtienen 0 puntos} \\
\hline
\end{tabularx}

\end{document}
